\chapter{Unternehmens- und Projektumfeld}
\label{cha:Unternehmens- und Projektumfeld}

\section{Vorstellung des Praxispartners}
Die vorliegende Projektarbeit wurde im Unternehmen MULTIVAC erarbeitet. MULTIVAC hat seinen Hauptsitz in Wolfertschwenden und besitzt eine weltweite Präsenz durch zahlreiche Tochtergesellschaften und Servicestandorte. Es ist ein familiengeführtes Unternehmen und wurde 1961 gegründet. MULTIVAC ist führender Anbieter von Verarbeitungs- und Verpackungslösungen für die Bereiche Lebensmittel, Medizin- und Pharmaprodukte sowie Konsum- und Industriegüter \autocite{UeberUns}. 

Das Produkt- und Leistungsportfolio des Unternehmens umfasst Thermoformverpackungsmaschinen und Traysealer, Kennzeichnungs- und Inspektionssysteme, Automatisierungslösungen sowie komplette Verpackungs- und Prozesslinien \autocite{Portfolio}.

In der Abteilung ,,Solutions'', in der die Projektarbeit realisiert wurde, werden kundenspezifische Produktions- und Verpackungslinien geplant und umgesetzt. In diesem Kontext werden verschiedene MULTIVAC-Maschinen und externe Komponenten zu einer durchgängigen Produktionslinie kombiniert. Die in dieser Arbeit erarbeitete Applikation soll hier das Testen beziehungsweise die Fehlersuche in der Maschinenkommunikation erleichtern.

\section{Maschinenkommunikation in Produktionslinien}
Die kundenspezifischen Produktionslinien sowie weitere Automatisierungslösungen setzen sich aus mehreren miteinander vernetzten Maschinen und Softwarekomponenten zusammen \autocite{Automatisierungsloesungen}. Die Gewährleistung eines reibungslosen Betriebs erfordert den Austausch von Zuständen, Ereignissen und Prozessdaten zwischen den Systemen. Die Kommunikation erfolgt dabei über definierte Schnittstellen und Protokolle \autocite[19--20]{EtherCAT}. 

Für echtzeitkritische Steuerungsaufgaben wird typischerweise das Feldbussysteme \textit{EtherCAT} eingesetzt. \textit{EtherCAT} überträgt zyklisch kleine Datenmengen mit fest definierter Struktur und garantiert sowohl kurze als auch deterministische Reaktionszeiten. Über diesen Kommunikationsweg werden beispielsweise Sensordaten, Aktorbefehle oder Zustandsinformationen zwischen Steuerung und Maschine ausgetauscht. \autocite[287--290]{EtherCAT}

Ergänzend dazu existiert eine übergeordnete Kommunikationsebene, auf welcher Informationen zwischen vollständigen Linienteilnehmern und der Liniensteuerung ausgetauscht werden. Hier kommt bei MULTIVAC das MQTT-Protokoll zum Einsatz. Im Gegensatz zu Feldbussen stehen hierbei weniger Echtzeitanforderungen im Vordergrund. Stattdessen werden größere und in ihrer Struktur variable Datenmengen übertragen. Beispiele hierfür sind Zustandsänderungen, Diagnosedaten oder Artikelinformationen \autocite{MQTTMessageQueue}.

Durch die höhere Informationsdichte sowie die Vielzahl unterschiedlicher Nachrichtentypen entsteht ein sehr komplexer Nachrichtenfluss. Während sich Feldbuskommunikation meist auf fest definierte Telegramme beschränkt, können MQTT-Nachrichten unterschiedlichste Inhalte besitzen \autocite[69]{IoTMQTT} und von zahlreichen Teilnehmern gleichzeitig veröffentlicht werden \autocite[11]{hivemqMQTTMQTT5}. Dies erschwert insbesondere bei komplexen Produktionslinien die Analyse von Kommunikationsabläufen und die systematische Fehlersuche.

\section{Herausforderungen bei der Analyse von Kommunikationsabläufen}
Bislang steht den Inbetriebnehmern und Software-Entwicklern lediglich eine Auswahl an Programmen zur Verfügung, die für die reine Aufzeichnung und Anzeige von Daten bestimmt sind. Die am häufigsten verwendeten Programme sind MQTT.fx und MQTT Explorer. Da diese für eine breite Nutzerbasis entwickelt wurden und allgemeine Anforderungen an die Analyse und Verwaltung von MQTT-Kommunikation abdecken, können unternehmens- und projektspezifische Anforderungen häufig nicht durch Standardfunktionen abgedeckt werden \autocite[13--14]{brooksNoSilverBullet1987}. Bei der Suche nach Fehlerursachen müssen daher häufig große Mengen an Nachrichten manuell durchsucht werden. Bleibt ein Fehler gänzlich unerkannt, kann das später schwerwiegende Folgen haben. Mögliche Fehler sind dabei beispielsweise fehlende Nachrichten zum Produktionsstand oder fehlerhafte Nachrichtenstrukturen, die andere Maschinen nicht verwerten können.

Ein weiterer Aspekt ist, dass eine Vielzahl von Inbetriebnehmern wenig Kenntnis hinsichtlich der Abläufe und des Nachrichtenaufbaus der Maschinenkommunikation verfügt. Infolgedessen sehen sich Software-Entwickler häufig veranlasst, die Maschine in der Produktionshalle aufzusuchen. Durch den Einsatz einer Analyse- und Monitoring-Applikation kann der erforderliche Zeitaufwand erheblich reduziert werden.